\subsection{Setting}

We study reward-based refinement of a pretrained generative policy.
An observation $s$ contains the inputs used by the policy, and an action
$a\in\mathbb{R}^{Hd}$ is a chunk of $H$ controls with $d$ dimensions each.
The frozen generator $\pi_0$ maps $s$ and noise $z$ to an action. Following
residual RL~\cite{dicerl2026}, we learn an additive correction:
\begin{equation}
  a_\theta(s,z)=\pi_0(s,z)+r_\theta(s,z),
  \qquad z\sim\mathcal{N}(0,I).
  \label{eq:residual}
\end{equation}
An ensemble critic estimates return from replayed experience. Our
manipulation tasks use sparse success rewards; DMC also includes dense
rewards. The goal is to improve task performance while retaining useful
pretrained behavior. Matched online comparisons share the base,
critic-training procedure, and interaction budget. Physical refinement
uses fixed experience collected before training.

\subsection{Accumulated Departure from the Base}

A critic can assign high values to actions outside its training coverage.
Repeatedly following those predictions may move the policy away from actions
that succeed~\cite{fujimoto2019bcq}. A step cap alone cannot prevent this:
$\|a_{t+1}-a_t\|_2\leq\delta$ permits distance from the initial action to
grow with the number of updates. Here $t$ indexes optimization steps.
Fig.~\ref{fig:problem} illustrates the issue by directly updating paired
actions under the same critic. In this toy, vanilla ascent moves probability
mass into an overestimated region, while a fixed-base projection retains
useful behavior nearby.

This distinction motivates two controls: a persistent restoring force
that opposes departure from the base, and a cap on each requested target
correction. We quantify departure using the residual's root-mean-square (rms) magnitude,
$D_{\mathrm{base}}=\mathbb{E}_{s,z}[\|r_\theta(s,z)\|_2/\sqrt{Hd}]$,
with states sampled from replay. The design question is how to resist
unnecessary changes while permitting useful ones. Because a neural actor
shares parameters across samples, constraints on its targets do not by
themselves guarantee equivalent constraints on its fitted outputs.
